\documentclass[12pt]{article}
\usepackage[T1]{fontenc}
\usepackage[utf8]{inputenc}
\usepackage{lmodern}
\usepackage{textcomp}
\usepackage{enumitem}
\usepackage{geometry}
\geometry{margin=1in}
\begin{document}
\begin{center}
Answer Key - Constitutional Law - Graduate Level Assessment 21 \
\small (Answers are ordered exactly as in the question paper.)
\end{center}

\section*{Multiple Choice}
\begin{enumerate}[label=\arabic*.]
\item \textbf{Answer:} B
\\ \textit{Options:} A) It is not a social rule.; B) It both confers power and imposes duties on judges to decide cases.; C) It is a primary rule, not a secondary rule.; D) It is applicable only in the case of federal constitutions.; E) It fails to specify the source of authority.
\\ \textit{Note:} The rule of recognition not only grants judges the authority to identify valid legal norms but also obligates them to apply these norms in their decisions, thereby embodying both power-conferring and duty-imposing characteristics.
\item \textbf{Answer:} A
\\ \textit{Options:} A) Integrity.; B) Hierarchy.; C) Dispute.; D) Justice.; E) Harmony.
\\ \textit{Note:} The correct answer is "Integrity" because Leopold Pospisil emphasizes that for a legal system to function effectively, it must not only possess authority and universality but also maintain a consistent and coherent application of laws, which is encapsulated in the notion of integrity.
\item \textbf{Answer:} D
\\ \textit{Options:} A) It is jurisdiction based on the country where the legal person was Registered; B) It is jurisdiction based on the nationality of the offender; C) It is jurisdiction based on the relationship between the victim and the offender; D) It is jurisdiction based on the nationality of the victims; E) It is jurisdiction based on the severity of the offence
\\ \textit{Note:} Passive personality jurisdiction refers to a legal principle that allows a state to exercise jurisdiction over a crime committed against its nationals, emphasizing the importance of protecting citizens regardless of where the crime occurs.
\item \textbf{Answer:} A
\\ \textit{Options:} A) The UK Constitution is a direct copy of the US Constitution; B) The UK Constitution is codified and contained in a single document; C) The UK Constitution can be changed by a simple majority vote in Parliament; D) The UK Constitution gives the judiciary the power to overturn acts of parliament; E) The UK Constitution is based on a Bill of Rights
\\ \textit{Note:} The statement is incorrect, as the UK Constitution is not a direct copy of the US Constitution; rather, it is an uncodified constitution comprising statutes, conventions, and legal precedents that differ significantly in structure and principles from the codified US Constitution.
\item \textbf{Answer:} A
\\ \textit{Options:} A) Vengeance.; B) Deterrence.; C) Restoration.; D) Moral education.; E) Indemnification.
\\ \textit{Note:} Durkheim viewed punishment primarily as a means of expressing societal outrage and reinforcing collective norms, which aligns with the concept of vengeance as it embodies the community's desire for retribution against violations of those norms.
\end{enumerate}

\section*{True / False}
\begin{enumerate}[label=\arabic*.]
\setcounter{enumi}{5}
\item \textbf{Answer:} True
\\ \textit{Options:} A) True; B) False
\\ \textit{Note:} Savigny's notion of Volksgeist posits that law is an embodiment of the historical, cultural, and social values of a community, emphasizing that legal systems arise from the collective consciousness of the people rather than being imposed by a singular authority.
\item \textbf{Answer:} False
\\ \textit{Options:} A) True; B) False
\\ \textit{Note:} The statement is false because Kelsen actually emphasized the importance of jurisprudence as a systematic study of law, rather than suggesting that legal norms overshadow the role of jurisprudence itself.
\item \textbf{Answer:} False
\\ \textit{Options:} A) True; B) False
\\ \textit{Note:} The statement is false because international law, particularly under the United Nations Charter, requires that the response to an armed attack must be proportional and necessary, meaning that the intensity of the initial armed force used is relevant in determining the appropriateness of a state's response.
\item \textbf{Answer:} False
\\ \textit{Options:} A) True; B) False
\\ \textit{Note:} The statement is false because the phrase 'Jurisprudence is the eye of law' is attributed to John Austin, not Jeremy Bentham, who is known for his utilitarian philosophy but did not use that specific phrase.
\item \textbf{Answer:} True
\\ \textit{Options:} A) True; B) False
\\ \textit{Note:} Weber's analysis of formally rational law highlights the role of rationalization and bureaucratic structures in legal development, which aligns with socialist principles that emphasize collective ownership and the systematic organization of society.
\end{enumerate}

\section*{Long Form Response}
\begin{enumerate}[label=\arabic*.]
\setcounter{enumi}{10}
\item \textbf{Answer:} A criminal statute prohibiting the dissemination of information on obtaining abortions raises significant constitutional concerns, particularly regarding First Amendment protections of commercial speech. Such a law infringes upon the right to publish information that is protected under the First Amendment, as it restricts the ability of individuals and organizations to communicate lawful services and options available to women. The Supreme Court has established that commercial speech, while subject to a lower level of protection than other forms of speech, is still safeguarded against outright bans unless justified by a substantial government interest that is narrowly tailored. In this case, the statute's broad prohibition does not adequately serve a compelling state interest without unduly limiting the free exchange of information essential to informed decision-making in healthcare, thus violating First Amendment rights.
\\ \textit{Note:} correct because it identifies the conflict between the criminal statute and the First Amendment's protection of commercial speech, emphasizing that such a law limits the dissemination of lawful information regarding abortion services, which is essential for informed decision-making.
\item \textbf{Answer:} Slander is legally defined as the oral communication of false statements that damage an individual's reputation. For a statement to qualify as slanderous, it is crucial that the accused intentionally disseminated the falsehood, demonstrating a degree of fault or intent. This intent can significantly influence the outcome of a legal case, as it distinguishes between mere negligence and willful defamation. If the speaker acted with malice or a reckless disregard for the truth, the implications are more severe, potentially leading to greater damages awarded to the victim. Ultimately, the requirement of intent underscores the balance between protecting free speech and safeguarding individuals from reputational harm.
\\ \textit{Note:} correct because it accurately defines slander as the oral spread of false statements harming reputation, while emphasizing the importance of intent in distinguishing between negligent and malicious actions, which is crucial for legal liability.
\end{enumerate}

\vfill
\noindent\textit{End of Answer Key}
\end{document}